\begin{lstlisting}[language=Java]

public class VarianceCalculator {

    // Calcula a variância pelo método tradicional.
    public static float calculateStandardVariance(int[] samples) {
        int sampleCount = samples.length;
        float mean = 0;

        // Calcula a média da amostra.
        for (int sample : samples) {
            mean += sample;
        }

        mean /= sampleCount;

        float squaredDifferenceSum = 0;

        // Soma os desvios ao quadrado em relação à média.
        for (int sample : samples) {
            float difference = sample - mean;
            squaredDifferenceSum += difference * difference;
        }

        return squaredDifferenceSum / sampleCount;
    }

    // Calcula a média da amostra de forma recursiva.
    private static float calculateWelfordMean(int[] samples, int index) {
        // Caso base: o primeiro elemento é a média dele mesmo.
        if (index == 0) {
            return samples[0];
        }

        float previousMean = calculateWelfordMean(samples, index - 1);

        // Atualiza a média considerando o novo elemento.
        return previousMean + (samples[index] - previousMean) / (index + 1);
    }

    // Soma recursivamente os desvios ao quadrado.
    private static float calculateSquaredDifferenceSum(int[] samples, int index, float mean) {
        float difference = samples[index] - mean;

        // Caso base da recursão.
        if (index == 0) {
            return difference * difference;
        }

        return calculateSquaredDifferenceSum(samples, index - 1, mean)
                + difference * difference;
    }

    // Calcula a variância utilizando os métodos recursivos.
    public static float calculateWelfordVariance(int[] samples) {
        int sampleCount = samples.length;

        float mean = calculateWelfordMean(samples, sampleCount - 1);

        float squaredDifferenceSum =
                calculateSquaredDifferenceSum(samples, sampleCount - 1, mean);

        return squaredDifferenceSum / sampleCount;
    }

    public static void main(String[] args) {
        int[] samples = {15, 29, 34, 40};

        System.out.println(calculateStandardVariance(samples)); // Saída 85.25
        System.out.println(calculateWelfordVariance(samples)); // Saída 85.25
    }
}

\end{lstlisting}